\documentclass[../../main.tex]{subfiles}

\graphicspath{{../../image/}{../../figure/}}

\begin{document}

\section{平衡律方程组的 $BV$ 解}\label{section:ch01-BV解}

此处考虑平衡律方程组\cref{align:ch01-closed-system} 的弱解 $U\in L^{\infty}(\mathscr{X})$，且 $U\in BV_{\mathrm{loc}}(\mathscr{X})$。此时，由\cref{theorem:ch01-BV链式法则}，函数 $G\circ U$ 也属于 $BV_{\mathrm{loc}}(\mathscr{X})\cap L^{\infty}(\mathscr{X})$，且\cref{align:ch01-closed-system} 作为测度等式成立。首要任务是结合\cref{theorem:ch01-BV-三分,theorem:ch01-BV链式法则,theorem:ch01-Gauss-Green} 考察\cref{align:ch01-closed-system} 的局部形式。

\begin{theorem}\label{theorem:ch01-BV解刻画}
函数 $U\in BV_{\mathrm{loc}}(\mathscr{X})\cap L^{\infty}(\mathscr{X})$ 是平衡律方程组\cref{align:ch01-closed-system} 的弱解，当且仅当：
\begin{enumerate}
\item\label{item:ch01-BV解刻画-a} 在 $U$ 的近似连续点集 $\mathscr{C}$ 上，测度等式
\begin{align}
[\mathsf{D}G(\tilde{U}(X),X),\operatorname{grad}U(X)]+\nabla\cdot G(\tilde{U}(X),X)=\Pi(\tilde{U}(X),X)\label{align:ch01-BV-continuous}
\end{align}
成立；
\item\label{item:ch01-BV解刻画-b} 跃度条件
\begin{align}
[G(U_{+},X)-G(U_{-},X)]N=0\label{align:ch01-BV-jump}
\end{align}
在 $U$ 的近似跳跃间断点集 $\mathscr{J}$ 上关于 $\mathscr{H}^{k-1}$ 几乎所有 $X$ 成立，其中 $N$ 为法向量、$U_{-},U_{+}$ 为单侧极限。
\end{enumerate}
\end{theorem}

\begin{proof}
在\cref{align:ch01-BV-continuous} 以及下文的\cref{align:ch01-BV-连续测度,align:ch01-BV-伴随测度} 中，符号 $\nabla\cdot$ 表示固定 $U$ 时关于 $X$ 的散度——区别于把 $U$ 视为 $X$ 函数的 $\operatorname{div}$。记 $\mathscr{M}$ 为以\cref{align:ch01-closed-system} 左端定义的测度。在 $\mathscr{C}$ 上，由\cref{theorem:ch01-BV链式法则}，$\mathscr{M}$ 化为\cref{align:ch01-BV-continuous} 左端的测度。回忆迹的\cref{definition:ch01-迹} 与\cref{theorem:ch01-BV-三分} 中单侧极限的刻画，我们推出在 $\mathscr{J}$ 的每一点有 $(G\circ U)_{\pm}=G\circ U_{\pm}$。于是若 $\mathscr{F}$ 为 $\mathscr{J}$ 的任一 Borel 子集，则由\cref{theorem:ch01-Gauss-Green} 证明后的注记，
\begin{align}
\mathscr{M}(\mathscr{F})=\int_{\mathscr{F}}[G(U_{-},X)-G(U_{+},X)]N\,\mathrm{d}\mathscr{H}^{k-1}.\label{align:ch01-BV-跳测度}
\end{align}
因此 $\mathscr{M}=\Pi$ 作为测度等式成立，当且仅当\cref{align:ch01-BV-continuous} 与\cref{align:ch01-BV-jump} 成立。证明完成。
\end{proof}

由此，$BV$ 解的近似跳跃间断点集是激波阵面的可数并。

如\cref{section:ch01-伴随平衡律} 中所见，当 $G$ 有伴随 $Q$ 时，伴随平衡律\cref{align:ch01-companion-law} 被\cref{align:ch01-closed-system} 的任一古典解自动满足。下述命题描述 $BV$ 弱解情形下的状况。

\begin{theorem}\label{theorem:ch01-BV伴随不等式}
设平衡律方程组\cref{align:ch01-closed-system} 配有伴随平衡律\cref{align:ch01-companion-law}。设 $U\in BV_{\mathrm{loc}}(\mathscr{X})\cap L^{\infty}(\mathscr{X})$ 为\cref{align:ch01-closed-system} 的弱解。则测度
\begin{align}
\mathscr{N}=\operatorname{div} Q(U(X),X)-h(U(X),X)\label{align:ch01-BV-伴随测度}
\end{align}
集中在 $U$ 的近似跳跃间断点集 $\mathscr{J}$ 上，且不等式\cref{align:ch01-admissibility} 在测度意义下成立，当且仅当
\begin{align}
[Q(U_{+},X)-Q(U_{-},X)]N\geqslant 0\label{align:ch01-BV-激波不等式}
\end{align}
对 $\mathscr{J}$ 上关于 $\mathscr{H}^{k-1}$ 几乎所有 $X$ 成立。
\end{theorem}

\begin{proof}
由\cref{theorem:ch01-BV链式法则}，可将\cref{align:ch01-closed-system} 与\cref{align:ch01-BV-伴随测度} 写为
\begin{align}
[\widetilde{\mathsf{D}G\circ U},\operatorname{grad}U]+\nabla\cdot G-\Pi=0,\label{align:ch01-BV-连续测度}
\end{align}
\begin{align}
\mathscr{N}=[\widetilde{\mathsf{D}Q\circ U},\operatorname{grad}U]+\nabla\cdot Q-h.\label{align:ch01-BV-伴随测度展开}
\end{align}

由\cref{align:ch01-normalized-composition}，若 $X$ 落在 $U$ 的近似连续点集 $\mathscr{C}$ 中，
\begin{align}
\widetilde{\mathsf{D}G\circ U}(X)=\mathsf{D}G(\tilde{U}(X),X),\qquad
\widetilde{\mathsf{D}Q\circ U}(X)=\mathsf{D}Q(\tilde{U}(X),X).\label{align:ch01-BV-规范化导数}
\end{align}
联立\cref{align:ch01-BV-连续测度,align:ch01-BV-伴随测度展开,align:ch01-BV-规范化导数}，并利用\cref{align:ch01-companion,align:ch01-companion-source}，我们推出 $\mathscr{N}$ 在 $\mathscr{C}$ 上消失。

由迹的\cref{definition:ch01-迹} 与\cref{theorem:ch01-BV-三分} 中单侧极限的刻画，我们推出 $(Q\circ U)_{\pm}=Q\circ U_{\pm}$。若 $\mathscr{F}$ 为 $\mathscr{J}$ 的有界 Borel 子集，则应用\cref{align:ch01-GG-jump}，同时记住\cref{theorem:ch01-Gauss-Green} 证明后的注记，得到
\begin{align}
\mathscr{N}(\mathscr{F})=\int_{\mathscr{F}}[Q(U_{-},X)-Q(U_{+},X)]N\,\mathrm{d}\mathscr{H}^{k-1}.\label{align:ch01-BV-伴随跳测度}
\end{align}
因此 $\mathscr{N}\leqslant 0$ 当且仅当\cref{align:ch01-BV-激波不等式} 成立。证明完成。
\end{proof}

\end{document}
